\subsection{实值函数的界}

定义 3. 设 \(f\) 是定义在集合 \(\mathbf{X}\) 上的实值函数，\(\mathbf{A}\) 是 \(\mathbf{X}\) 的一个非空子集。此时集 \(f(\mathbf{A})\) 在 \(\overline{\mathbf{R}}\) 中的上确界（相应地，下确界）称为 \(f\) 在 \(\mathbf{A}\) 中的上确界（相应地，下确界），记作 \(\sup_{x\in \mathbf{A}}f(x)\left[\text{resp. } \inf_{x\in \mathbf{A}} f(x)\right]\) 。

特别地，若 A 是 \(\overline{R}\) 的一个非空子集，则

\[
\sup \mathrm{A} = \sup _ {x \in \mathrm{A}} x.\tag{1}
\]

用右端的记号来表示 A 的上确界，往往更为方便。

数 \(a = \sup_{x \in \mathbf{A}} f(x)\) 由下列两条性质刻画：

\begin{enumerate}
\def\labelenumi{(\roman{enumi})}
\item
  对一切 \(x \in \mathbf{A}, f(x) \leqslant a\) 。
\item
  对每个 \(b < a\) ，存在 \(x \in \mathbf{A}\) 使得 \(b < f(x) \leqslant a\) 。
\end{enumerate}

数 \(\sup_{x\in\mathbf{A}}f(x)\) 与 \(\inf_{x\in\mathbf{A}}f(x)\) 属于 \(f(\mathbf{A})\) 在 \(\overline{\mathbf{R}}\) 中的闭包。我们有 \(\inf_{x\in\mathbf{A}}f(x)\leqslant\sup_{x\in\mathbf{A}}f(x)\) ；且这两个数相等当且仅当 \(f\) 在 \(\mathbf{A}\) 上取常值。

定义在集合 X 上的实值函数在 X 的一个非空子集 A 中上有界（相应地，下有界）当且仅当 \(\sup_{x\in\mathbf{A}}f(x)<+\infty\) {[}相应地，\(\inf_{x\in\mathbf{A}}f(x)>-\infty\) {]}。f 在 A 中有界当且仅当 \(|f|\) 在 A 中上有界，亦即当且仅当 \(\sup_{x\in\mathbf{A}}|f(x)|<+\infty\) 。

我们有

\[
\inf _ {x \in \mathbf {A}} f (x) = - \sup _ {x \in \mathbf {A}} (- f (x)).\tag{2}
\]

这一关系式把下确界的所有性质都化归为上确界的性质；因此在一般情况下我们只谈论后者。

命题 5. 设 \(f\) 是定义在集合 \(\mathbf{X}\) 上的一个实值函数。在集 \(\mathfrak{F}(\mathbf{X})\) （由 \(\mathbf{X}\) 的所有有限子集组成，按关系 \(\subset\) 有向）上，实值函数 \(\mathbf{H} \to \sup_{x \in \mathbf{H}} f(x)\) 是递增的，实值函数 \(\mathbf{H} \to \inf_{x \in \mathbf{H}} f(x)\) 是递减的，并且我们有

\[
\left\{ \begin{array}{l} \sup _ {x \in \mathbf {A}} f (x) = \lim _ {\mathbf {H} \in \mathfrak {F} (\mathbf {X})} (\sup _ {x \in \mathbf {H}} f (x)), \\ \inf _ {x \in \mathbf {A}} f (x) = \lim _ {\mathbf {H} \in \mathfrak {F} (\mathbf {X})} (\inf _ {x \in \mathbf {H}} f (x)). \end{array} \right.\tag{3}
\]

设 \(\varphi(\mathbf{H}) = \sup_{x \in \mathbf{H}} f(x)\) 。显然 \(\varphi\) 递增，因而有极限 \(a\) （第 2 节，定理 2）；又因对一切 H 有 \(\varphi(\mathbf{H}) \leqslant \sup_{x \in \mathbf{A}} f(x)\) ，故有 \(a \leqslant \sup_{x \in \mathbf{A}} f(x)\) （第 2 节，定理 1）。倘若 \(a < \sup_{x \in \mathbf{A}} f(x)\) ，就会存在 \(x_0 \in \mathbf{X}\) 使得 \(a < f(x_0)\) ；但这是矛盾的，因为 \(\varphi(\mathbf{H}) \geqslant f(x_0)\) （只要 \(x_0 \in \mathbf{H}\) ）。

特别地，由 (1)，若 A 是 \(\overline{R}\) 的任意一个非空子集，则我们有

\[
\sup \mathrm{A} = \lim _ {\mathrm{H} \in \mathfrak {F} (\mathrm{A})} (\sup _ {x \in \mathrm{H}} x).\tag{4}
\]

命题 6. 设 \(f\) 与 \(g\) 是定义在 \(\mathbf{X}\) 上的两个实值函数。若 \(f(x) \leqslant g(x)\) 对每一点 \(x\) （属于某个非空子集 \(A\) ，而该子集含于 \(\mathbf{X}\) ）成立，则我们有

\[
\left\{ \begin{array}{l} \sup _ {x \in \mathbf {A}} f (x) \leqslant \sup _ {x \in \mathbf {A}} g (x), \\ \inf _ {x \in \mathbf {A}} f (x) \leqslant \inf _ {x \in \mathbf {A}} g (x). \end{array} \right.\tag{5}
\]

命题 7. 设 \(f\) 是定义在 \(\mathbf{X}\) 上的一个实值函数。若 \(\mathbf{A}\) 与 \(\mathbf{B}\) 是 \(\mathbf{X}\) 的两个非空子集，且 \(\mathbf{A} \subset \mathbf{B}\) ，则

\[
\sup _ {x \in \mathbf {A}} f (x) \leqslant \sup _ {x \in \mathbf {B}} f (x).\tag{6}
\]

命题 8. 设 \(f\) 是定义在 \(\mathbf{X}\) 上的一个实值函数，且设 \((\mathbf{A}_t)_{t \in I}\) 是由 \(\mathbf{X}\) 的非空子集组成的一个非空族；则

\[
\sup_{x\in \bigcup_{t\in \mathbf{I}}\mathbf{A}_{t}}f(x) = \sup_{t\in \mathbf{I}}(\sup_{x\in \mathbf{A}_{t}}f(x)).\tag{7}
\]

设 \(f\) 是定义在乘积集合 \(\mathbf{X}_1 \times \mathbf{X}_2\) 上的一个实值函数。若 \(\mathbf{A}_2\) 是 \(\mathbf{X}_2\) 的一个非空子集，我们用 \(\sup_{x_2 \in \mathbf{A}_2} f(x_1, x_2)\) 表示在 \(\mathbf{A}_2\) 中取的上确界，即实值函数 \(x_2 \to f(x_1, x_2)\) （定义在 \(\mathbf{X}_2\) 上）的上确界。由命题 8 我们特别推出：

命题 9. 设 \(f\) 是定义在乘积集合 \(\mathbf{X}_1 \times \mathbf{X}_2\) 上的一个实值函数。若 \(\mathbf{A}_1, \mathbf{A}_2\) 分别是 \(\mathbf{X}_1, \mathbf{X}_2\) 的任意非空子集，则

\[
\sup _ {(x _ {1}, x _ {2}) \in \mathbf {A} _ {1} \times \mathbf {A} _ {2}} f (x _ {1}, x _ {2}) = \sup _ {x _ {1} \in \mathbf {A} _ {1}} (\sup _ {x _ {2} \in \mathbf {A} _ {2}} f (x _ {1}, x _ {2})) = \sup _ {x _ {2} \in \mathbf {A} _ {2}} (\sup _ {x _ {1} \in \mathbf {A} _ {1}} f (x _ {1}, x _ {2})).\tag{8}
\]

